\chapter{Problema Abstracto de Cauchy}
\label{ch:cauchy}
\todo[inline]{HABLAR DE LOS TIEMPOS, PORQUE A VECES ES $t_0 < t < t_1$ Y A VECES $0 < t_0 < t$}
La utilidad de la teoría de semigrupos de operadores lineales radica, entre otras cosas, en su utilidad como lenguaje para estudiar la dinámica en espacios funcionales; en concreto, en espacios de Banach. 

A primera vista, esto pudiese parecer una cuestión limitada a la matemática abstracta, e incluso uno podría llegar a pensar que es necesario hacer un esfuerzo para concebir problemas de evolución típicos en las ciencias aplicadas como una dinámica en un espacio funcional. La realidad no podría ser más distinta, y es que, en realidad, muchas de las ecuaciones en derivadas parciales (que se encuentran, como disciplina, en un lugar cómodo entre la difusa línea de la matemática aplicada y fundamental) que se conocen se pueden plantear como una ecuación de evolución en un espacio funcional. 

Veamos un ejemplo. Consideremos el problema de valor inicial y de frontera asociado a la ecuación de onda:

\begin{align}
\label{eqOnda}
    \begin{cases}
        u_{tt}(x,t) = c^2 \nabla^2 u(x,t), \ &(x,t) \in \Omega \times \mathbb{R}^+,\\
        u(x,0) = u_0(x), \ &x \in \Omega,\\
        u_t(x,0)  = u_1(x) \ &x \in \Omega,\\
        u(x,t) = 0 \ &(x,t) \in \partial \Omega \times \mathbb{R}^+ 
    \end{cases}
\end{align}

Si consideramos el cambio de variable:

\[v(x,t) = u_t(x,t),\]

entonces el problema anterior se puede reescribir como el sistema:

\begin{align}
\label{eqOndaCambioDeVar}
    \begin{cases}
        u_{tt}(x,t) = c^2 \nabla^2 u(x,t), \ &(x,t) \in \Omega \times \mathbb{R}^+,\\
        v(x,t) = u_t(x,t), \ &(x,t) \in \Omega \times \mathbb{R}^+,\\
        u(x,0) = u_0(x), \ &x \in \Omega,\\
        v(x,0)  = v_0(x) \ &x \in \Omega,\\
        u(x,t) = 0 \ &(x,t) \in \partial \Omega \times \mathbb{R}^+,
    \end{cases}
\end{align}

donde (por una argumentación que ahora mismo no es constructiva para nuestro ejemplo) se tiene que:

\[u(\cdot, t) \in H_0^1(\Omega), \ v(\cdot, t) \in L^2(\Omega).\]

Ahora definimos:

\begin{align*}
    U(t) &= \big(u(\cdot, t), v(\cdot, t)\big) \in H_0^1(\Omega) \times L^2(\Omega)\\
    U(0) &= U_0
\end{align*}

y tenemos que el sistema anterior se puede reescribir en términos de $U$ y un operador lineal $A$ con

\[A = 
    \begin{pmatrix}
    0 & I \\
    c^2\nabla^2 & 0
    \end{pmatrix},
\]

como un sistema dinámico en un espacio funcional. Entiéndase:

\begin{align}
    \label{eqOndaFuncional}
    U_t &= AU, \quad t \in \mathbb{R}^+,\\
    U(0) &= U_0 \in H_0^1(\Omega) \times L^2(\Omega).
\end{align}

Vemos entonces que, mediante simple manipulación, hemos sido capaces de reescribir uno de los problemas más clásicos en las ciencias aplicadas (como lo es la ecuación de onda) en un problema de evolución en un espacio funcional. La utilidad de la teoría de semigrupos se presentará cuando el operador $A$ del sistema anterior (y, en general, de cualquier sistema con dicha forma) sea el generador infinitesimal de un semigrupo de operadores lineales. Cuando ese sea el caso, tendremos derecho a usar todos los resultados discutidos en capítulos anteriores que, entre otras cosas, nos permitirán demostrar la existencia y unicidad de soluciones a estos problemas de evolución.

Dicho todo esto y habiendo entendido el uso que haremos de la teoría de semigrupos, demos inicio a la primera sección de este capítulo formalizando qué entendemos por un \textit{problema de Cauchy homogéneo} en un espacio de Banach, cómo definimos sus soluciones y un breve estudio de cómo podemos probar la existencia y unicidad de las mismas.

\section{Introducción: Problema de Cauchy homogéneo}
\label{ch:cauchy:section:Hom}

En primer lugar, recordemos que un problema de Cauchy es, en teoría de ecuaciones diferenciales ordinarias, un problema de valor inicial. Así, pues, un problema de Cauchy abstracto (o, simplemente, problema de Cauchy) homogéneo en un espacio de Banach $E_0$ tiene la forma:

\begin{align}
    \label{ch:cauchy:section:Hom:eqBasica}
    \frac{de}{dt} &= Ae, \quad 0 \leq t_0 < t\\
    e(t_0) &= e_0 \in E_0,
\end{align}

donde $A$ es el generador infinitesimal de un semigrupo fuertemente continuo $\{T(t) ; t \geq 0\} \subset L(E_0)$.

Damos ahora la definición de una solución del problema \eqref{ch:cauchy:section:Hom:eqBasica}. En realidad, podremos distinguir dos tipos de soluciones, una \textit{solución fuerte} y una \textit{solución débil}:\\

\begin{Definicion}{(Soluciones al problema de Cauchy homogéneo)}
    \label{ch:cauchy:section:Hom:def:soluciones}

    \begin{enumerate}[label=\alph*)]
        \item Se define una \textit{solución fuerte} de \eqref{ch:cauchy:section:Hom:eqBasica} como una función $e : [0,\infty) \rightarrow E_0$ tal que verifica \eqref{ch:cauchy:section:Hom:eqBasica}, $e(t_0) = e_0$, $e$ es continuamente diferenciable y $e(t) \in D(A)$ para $t > 0$.

        \item Por otro lado, \textit{solución débil} de \eqref{ch:cauchy:section:Hom:eqBasica} es una función continua $e : [0,\infty) \rightarrow E_0$ tal que $e(t_0) = e_0$ y para todo $e^* \in D(A^*)$ se verifica que

        \[\frac{d^+}{dt}\langle e^*, e(t) \rangle = \langle A^*e^*, e(t)\rangle.\]
    \end{enumerate}
\end{Definicion}

La definición anterior nos revela casi trivialmente que una solución fuerte es también una solución débil. Veamos ahora la forma que podría tomar una de estas soluciones débiles.

Para ello, consideramos la aplicación

\[e(t) = T(t-t_0)e_0,\]

donde $e_0 \in E_0$ y $T(t)$ es el semigrupo del cual $A$ es generador. Comprobemos ahora que $e(t)$ es una solución débil. Para ello, tomamos $e^* \in E_0^*$ y desarrollamos $\displaystyle\frac{d}{dt}\langle e^*, e(t)\rangle$:

\begin{align*}
    \frac{d}{dt}\langle e^*, e(t)\rangle &= \frac{d}{dt}\langle e^*, T(t-t_0)e_0\rangle\\
    &= \frac{d}{dt}\Big\langle \big(T(t-t_0)\big)^*e^*, e_0\Big\rangle,
\end{align*}

donde \textcolor{red}{SEMIGRUPO DUAL...} y, como $e^* \in D(A^*)$, entonces

\begin{align*}
    \frac{d}{dt}\langle e^*, e(t)\rangle &= \frac{d}{dt}\Big\langle \big(T(t-t_0)\big)^*e^*, e_0\Big\rangle\\
    &= \left \langle\frac{d}{dt}\big(T(t-t_0)\big)^*e^*, e_0 \right\rangle\\
    &= \Big\langle A^*\big(T(t-t_0)\big)^*e^*, e_0\Big\rangle\\
    &= \Big\langle \big(T(t-t_0)\big)^*A^*e^*, e_0\Big\rangle\\
    &=\langle A^*e^*, T(t-t_0)e_0\rangle\\
    &=\langle A^*e^*, e(t) \rangle.
\end{align*}

Por lo tanto, $e(t)$ es una solución débil. Veamos ahora que esta solución es única, probando así que todas las soluciones débiles (y, en consecuencia, las fuertes pues estas también son débiles) son de la forma $e(t) = T(t-t_0)e_0$. Para ello, sean $e_1(t)$ y $e_2(t)$ dos soluciones débiles al problema homogéneo, entonces

\[u(t) = e_1(t) - e_2(t)\]

es también una solución al problema homogéneo con dato inicial

\[u(t_0) = e_1(t_0) - e(t_1) = e_0 - e_0 = 0,\]

pues

\begin{align*}
    \frac{d^+}{dt}\langle e^*, u(t) \rangle &= \frac{d^+}{dt}\langle e^*, e_1(t)\rangle - \frac{d^+}{dt}\langle e^*, e_2(t)\rangle\\
    &= \langle e^* Ae_1(t)\rangle - \langle e^* Ae_2(t)\rangle\\
    &= \langle e^*, u(t) \rangle.
\end{align*}

Construyamos ahora una función con más regularidad. Para ello, definimos

\[U(t) = \int_{t_0}^t u(s) \, ds,\]

donde $U$ es de clase $C^1$ y claramente

\[\frac{d^+}{dt}U(t) = u(t).\]

Entonces:

\begin{align*}
    \langle e^*, u(t) \rangle &= \int_{t_0}^t \frac{d^+}{dt}\langle e^*, u(s)\rangle \, ds\\
    &= \int_{t_0}^t \langle e^*, Au(s)\rangle \, ds\\
    &= \int_{t_0}^t \langle A^*e^*, u(s)\rangle \, ds\\
    &= \left \langle A^*e^*, \int_{t_0}^t u(s) \, ds\right \rangle\\
    &= \langle A^*e^*, U(t) \rangle,
\end{align*}

y así

\[\langle e^*, u(t) \rangle = \left \langle e^*, \frac{d^+}{dt}U(t)\right \rangle = \langle A^*e^*, U(t) \rangle.\]

Construimos ahora la función auxiliar

\[t \mapsto \langle e^*, T(t^* - t)U(t) \rangle,\]

donde $t_0 \leq t \leq t^*$ y $t^*$ es un instante arbitrario. Ahora derivamos la función anterior y obtenemos:

\begin{align*}
    \frac{d^+}{dt}\langle e^*, T(t^* - t)U(t) \rangle &= \frac{d^+}{dt}\Big \langle \big (T(t^* - t)\big)^*e^*, U(t) \Big\rangle\\
    &= \left \langle \frac{d^+}{dt} \big (T(t^* - t)\big)^*e^*, U(t) \right\rangle + \left\langle \big (T(t^* - t)\big)^*e^*, \frac{d^+}{dt}U(t) \right\rangle\\
    &= -\Big\langle A^*\big (T(t^* - t)\big)^*e^*, U(t) \Big\rangle + \left\langle \big (T(t^* - t)\big)^*e^*, \frac{d^+}{dt}U(t) \right\rangle.
\end{align*}

Para continuar, queremos usar sobre el término

\[\left\langle \big (T(t^* - t)\big)^*e^*, \frac{d^+}{dt}U(t) \right\rangle\]

la igualdad

\[\left \langle e^*, \frac{d^+}{dt}U(t)\right \rangle = \langle A^*e^*, U(t) \rangle.\]

Para ello, hay que probar que $\big(T(t^*-t)\big)^*e^*\in D(A^*)$, donde $e^* \in D(A^*)$. Tomando $e \in D(A)$:

\[\Big\langle \big(T(t^*-t)\big)^*e^*, Ae \Big\rangle = \langle e^*, T(t^* - t)Ae \rangle = \langle A^*e^*, T(t^*-t)e\rangle,\]

y así, por definición de $D(A^*)$, se tiene que $\big(T(t^*-t)\big)^*e^*\in D(A^*)$ para cualquier $e^* \in D(A^*)$.

Dicho eso, aplicamos la igualdad mencionada y obtenemos:

\begin{align*}
    \frac{d^+}{dt}\langle e^*, T(t^* - t)U(t) \rangle &= -\Big\langle A^*\big (T(t^* - t)\big)^*e^*, U(t) \Big\rangle + \left\langle \big (T(t^* - t)\big)^*e^*, \frac{d^+}{dt}U(t) \right\rangle\\
    &= -\Big\langle A^*\big (T(t^* - t)\big)^*e^*, U(t) \Big\rangle + \Big\langle A^*\big (T(t^* - t)\big)^*e^*, U(t) \Big\rangle\\
    &= 0,
\end{align*}

por lo que la función $t \mapsto \langle e^*, T(t^* - t)U(t) \rangle$ es constante y, como $U(t_0) = 0$, entonces es nula. Evaluándola en $t^* > t_0$ arbitrario:

\[\langle e^*, T(t^* - t^*)U(t^*)\rangle = \langle e^*, U(t^*) \rangle = 0 \ ; \ t^* \in (t_0, \infty).\]

Por ser cierto para todo $e^* \in D(A^*)$, lo anterior implica que $U(t^*) = 0$ para todo $t^* \in [t_0, \infty)$ (hace falta una sutileza argumentativa en este punto que se puede encontrar en la prueba del Teorema \eqref{ch:cauchy:section:noHom:th:solDebilYFuerteRels}; sin embargo, dado que estamos en una introducción, nos permitimos seguir sin ella). Así, $u(s) = 0$ para todo $s \in [t_0, \infty)$, lo cual prueba la unicidad de soluciones.

De todo lo anterior se obtiene entonces que, de existir, tanto las soluciones débiles como las fuertes del problema \eqref{ch:cauchy:section:Hom:eqBasica} son de la forma 

\begin{align}
    \label{Hom:FormaSoluciones}
    e(t) = T(t-t_0)e_0,
\end{align}

donde $T(t)$ es el semigrupo fuertemente continuo generado por $A$ y $e_0$ el dato inicial del problema homogéneo. 

Nos hacemos ahora la siguiente pregunta sobre la regularidad de las soluciones débiles: ¿Bajo qué condiciones una solución débil pasa a ser una solución fuerte? Veamos que una solución débil es, en realidad, una solución fuerte cuando el dato inicial $e_0$ pertenece a $D(A)$.

Para ello, debemos comprobar que si $e$ es una solución débil al problema, entonces $e(t) \in D(A)$ para todo instante $t$, $e$ satisface \eqref{ch:cauchy:section:Hom:eqBasica} y $e$ es continuamente diferenciable.

Que $e(t) \in D(A)$ para todo instante $t$ es trivial, pues sabemos por \eqref{Hom:FormaSoluciones} que $e(t) = T(t-t_0)e_0$, y como $e_0 \in D(A)$ entonces se tiene el resultado. Por otro lado, $t \mapsto e(t)$ es continua por continuidad del semigrupo $T(t)$. Veamos ahora que $e$ satisface \eqref{ch:cauchy:section:Hom:eqBasica}:

\begin{align*}
    \frac{d^+}{dt}e(t) = \frac{d^+}{dt}T(t-t_0)e_0 = AT(t-t_0)e_0 = Ae(t),
\end{align*}

por lo que $e$ satisface \eqref{ch:cauchy:section:Hom:eqBasica}. Por último, debemos comprobar que la derivada de $e$ es continua. Para ello, usamos la conmutatividad del semigrupo con el generador y obtenemos:

\[\frac{d^+}{dt}e(t) = AT(t-t_0)e_0 = T(t-t_0)Ae_0,\]

y como $Ae_0 \in E_0$ es un elemento del espacio cualquiera, entonces volvemos a invocar la continuidad del semigrupo y se tiene inexorablemente que la derivada es continua. En consecuencia, una solución débil con un dato inicial regular dan como resultado una solución fuerte.

Habiendo definido qué entendemos por soluciones del problema \eqref{ch:cauchy:section:Hom:eqBasica} y qué formas pueden tener dichas soluciones, damos fin a esta sección estudiando bajo qué condiciones podemos demostrar la existencia y unicidad de soluciones fuertes, pues la de las soluciones débiles se desprende (como ya vimos anteriormente) del hecho de que $A$ sea el generador infinitesimal de un semigrupo fuertemente continuo; es decir, el problema homogéneo \eqref{ch:cauchy:section:Hom:eqBasica} posee soluciones débiles si y sólo si $A$ es el generador infinitesimal de un semigrupo fuertemente continuo.

Exigiendo menos propiedades del operador $A$, en \textcolor{green}{[PAZY CAPÍTULO 4.1]} se comenta cómo podemos obtener soluciones incluso en ecuaciones de evolución de la forma \eqref{ch:cauchy:section:Hom:eqBasica} donde $A$ no es el generador infinitesimal de un semigrupo fuertemente continuo; sin embargo, y continuando con la cita al texto referenciado, para obtener la existencia y unicidad de soluciones diferenciables en $[0, \infty)$ cuyo dato inicial $e_0$ pertenezca a $D(A)$, es necesario asumir que $A$ es el generador infinitesimal de un semigrupo fuertemente continuo. Esto, junto con todo el estudio anterior, da pie al siguiente teorema:\\

\begin{Teorema}{(Existencia y unicidad de soluciones del problema de Cauchy homogéneo)}
    \label{cauchy:hom:ExistenciaYUnicidadFuerte}

    Sea $A$ un operador lineal densamente definido con conjunto resolvente $\rho(A)$ no vacío. El problema de valor inicial \eqref{ch:cauchy:section:Hom:eqBasica} tiene, para cada dato inicial $e_0 \in D(A)$, una única solución $e(t)$ que es continuamente diferenciable en $[0, \infty)$ si y sólo si $A$ es el generador infinitesimal de un semigrupo fuertemente continuo.
\end{Teorema}

\begin{proof}
    Si asumimos que $A$ es el generador infinitesimal de un semigrupo fuertemente continuo, entonces la prueba de la implicación es, esencialmente, todo la discusión sostenida anteriormente. La otra implicación se puede encontrar en \textcolor{green}{[PAZY CAPÍTULO 4.1]} y la omitimos pues no resulta demasiado relevante para este texto.
\end{proof}

Como se podrá ver en lo sucesivo, durante este capítulo estaremos repitiendo una estructura clara entre secciones: definición de los tipos de soluciones, forma de las soluciones débiles y resultados o bien de regularidad o de existencia y unicidad. 

\section{Problema de Cauchy no homogéneo}
\label{ch:cauchy:section:noHom}

En esta sección estudiaremos problemas de Cauchy (es decir, problemas de valor inicial) en un espacio de Banach $E_0$ para ecuaciones lineales no homogéneas de la forma:\textcolor{red}{ARREGLAR EL LABEL}

\begin{align}
    \label{ch:cauchy:section:noHom:eqBasica}
    \frac{de}{dt} &= Ae + f(t), \quad t_0 < t < t_1\\
    e(t_0) &= e_0 \in E_0,
\end{align}
\footnotetext{Esto es exclusivamente notación salvo que se indique lo contrario.}
donde $A$ es el generador infinitesimal de un semigrupo fuertemente continuo $\{\mathrm{e}^{At} ; t \geq 0\} \subset L(E_0)$\footnotemark y $f : [t_0, t_1) \rightarrow E_0$ es una función continua por partes y continua por la derecha.

Empezamos entonces por definir los tipos de soluciones que puede tener el problema \eqref{ch:cauchy:section:noHom:eqBasica}\\

\begin{Definicion}{(Soluciones al problema de Cauchy no homogéneo)}
    \label{ch:cauchy:section:noHom:def:soluciones}

    \begin{enumerate}[label=\alph*)]
        \item Una \textit{solución fuerte} de \eqref{ch:cauchy:section:noHom:eqBasica} es una función continua $e : [t_0, t_1) \rightarrow E_0$ tal que $e(t_0) = e_0$ y para $t_0 < t < t_1$ se verifica que $e(t) \in D(A)$, 

        \begin{align}
        \label{ch:cauchy:section:noHom:def:soluciones:1}
        \frac{d^+}{dt}e(t) = \lim_{h \rightarrow 0^+} \frac{e(t+h) - e(t)}{h}
        \end{align}

        existe, se satisface \eqref{ch:cauchy:section:noHom:eqBasica} con $\displaystyle\frac{d}{dt}e(t)$ sustituida por $\displaystyle\frac{d^+}{dt}e(t)$ y $t \mapsto \frac{d^+}{dt}e(t)$ es continua donde lo sea $f$. Nótese que se $t \mapsto f(t)$ es una función continua y $t \mapsto e(t)$ es una solución fuerte de \eqref{ch:cauchy:section:noHom:eqBasica}, entonces $t \mapsto e(t)$ es continuamente diferenciable y \eqref{ch:cauchy:section:noHom:eqBasica} se verifica para cada $t \in (t_0, t_1)$.

        \item Una \textit{solución débil} de \eqref{ch:cauchy:section:noHom:eqBasica} en $[t_0, t_1)$ es una función continua $e : [t_0, t_1) \rightarrow E_0$ tal que $e(t_0) = e_0$ y para todo $e^* \in D(A^*)$, $t\mapsto \langle e^*, e(t) \rangle$ tiene derivada por la derecha y

        \begin{align}
        \label{ch:cauchy:section:noHom:def:soluciones:2}
        \frac{d^+}{dt}\langle e^*, e(t)\rangle = \langle A^* e^*, e(t) \rangle + \langle e^*, f(t) \rangle, \quad t_0 < t < t_1.
        \end{align}

        Nótese que si $t \mapsto \langle e^*, f(t) \rangle$ es una función continua y $t \mapsto e(t)$ es una solución débil de \eqref{ch:cauchy:section:noHom:eqBasica}, entonces $t \mapsto \langle e^*, e(t) \rangle$ es continuamente diferenciable y \eqref{ch:cauchy:section:noHom:def:soluciones:2} se verifica con $\displaystyle\frac{d^+}{dt}$ sustituido por $\displaystyle\frac{d}{dt}$
    \end{enumerate}
\end{Definicion}

Veamos ahora dos definiciones que nos serán útiles tanto en el estudio de la existencia y unicidad de soluciones del problema \eqref{ch:cauchy:section:noHom:eqBasica}, como en el estudio de la densidad de dominios.\\

\begin{Definicion}{(Subconjunto total y anulador)}
    \label{ch:cauchy:section:noHom:def:totalYAnulador}
    Un subconjunto $S^* \subset E^*_0$ se dice \textit{total} si

    \[e \in E_0, \langle e^*, e \rangle = 0 \ \  \forall e^* \in S^* \implies e = 0.\]

    Por otro lado, el \textit{anulador} $S^\perp \subset E^*_0$ de un subconjunto $S \subset E_0$ es el conjunto de todos los elementos $e^* \in E_0^*$ tales que $\langle e^*, e \rangle = 0 \ \forall e \in S$. Sabemos que $S \subset E_0$ es un espacio vectorial, por lo que $(S^\perp)^\perp = \overline{S}$\textcolor{red}{DEMOSTRAR} \textcolor{green}{(EL LIBRO INDICA H. Br´ezis. Analyse fonctionnelle)}
\end{Definicion}

Continuamos ahora con un Lema que nos permitirá clasificar el dominio del adjunto de un operador como total.\\

\begin{Lema}{(Condición suficiente para que el dominio del adjunto sea total)}
    \label{ch:cauchy:section:noHom:lema:dominioAdjTotal}

    Si $A : D(A) \subset E_0 \rightarrow E_0$ es cerrado y densamente definido, entonces $D(A^*)$ es total.
\end{Lema}

\begin{proof}
    Sea $e \in E_0$ tal que $\langle e^*, e \rangle = 0$ para todo $e^* \in D(A^*)$. Queremos demostrar que $e = 0$. Nótese que el gráfico de $A^*$, $G(A^*) = \{(e^*, A^*e^*) : e^* \in D(A^*)\}$ es el anulador en $E_0^* \times E_0^*$ de $\Gamma = \{(-Ae, e) : e \in D(A) \}$ (pues $\langle (e^*, A^*e^*), (-Ae, e)\rangle = \langle e^*, -Ae\rangle + \langle A^*e^*, e\rangle = -\langle e^*, Ae\rangle + \langle e^*, Ae\rangle = 0$); esto es, $\Gamma^\perp = G(A^*)$. Como $G(A^*)$ también anula a $(e,0)$ (porque $\langle (e^*, A^*e^*), (e,0)\rangle = \langle e^*, e\rangle + \langle A^*e^*, 0 \rangle = 0$), entonces se tiene que $(e,0) \in G(A^*)^\perp = \overline{\Gamma} = \Gamma$ (esa última igualdad se deduce del hecho de que $\Gamma$ es cerrado por serlo $A$), por lo que existe $x \in D(A)$ tal que $(e,0) = (-Ax, x) \implies x = 0, e = -A0 = 0$. Es decir, $e = 0$ tal y como queríamos probar.
\end{proof}

Adentrándonos más en el terreno de las soluciones al problema de Cauchy abstracto planteado \eqref{ch:cauchy:section:noHom:eqBasica}, se presenta ahora un teorema que nos permitirá manipular de manera más simple las soluciones débiles (literalmente otorgándonos la forma que deben tener), así como establecer algunas relaciones entre soluciones fuertes y débiles.\\

\begin{Teorema}{(Manipulación de soluciones débiles y relaciones entre soluciones débiles y fuertes)}
    \label{ch:cauchy:section:noHom:th:solDebilYFuerteRels}

    \begin{enumerate}
        \item Si $e : [t_0, t_1) \rightarrow E_0$ es una solución fuerte de \eqref{ch:cauchy:section:noHom:eqBasica}, entonces es también una solución débil de \eqref{ch:cauchy:section:noHom:eqBasica}.

        \item Si $e : [t_0, t_1) \rightarrow E_0$ es una solución débil de \eqref{ch:cauchy:section:noHom:eqBasica}, entonces

        \begin{align}
            \label{ch:cauchy:section:noHom:th:solDebilYFuerteRels:2}
            e(t) = \mathrm{e}^{A(t-t_0)}e_0 + \int_{t_0}^{t}\mathrm{e}^{A(t-s)}f(s) \, ds, \quad t_0 \leq t \leq t_1
        \end{align}

        En particular, existe una única solución débil de \eqref{ch:cauchy:section:noHom:eqBasica}.

        \item Si $e : [t_0, t_1) \rightarrow E_0$ está definido por \eqref{ch:cauchy:section:noHom:th:solDebilYFuerteRels:2}, entonces $e : [t_0, t_1) \rightarrow E_0$ es una solución débil de \eqref{ch:cauchy:section:noHom:eqBasica}.

        \item Sea $e : [t_0, t_1) \rightarrow E_0$ una solución débil de \eqref{ch:cauchy:section:noHom:eqBasica}. Si para algún instante $t \in (t_0, t_1)$ se cumple que $e(t) \in D(A)$ o bien que $\displaystyle\frac{d^+}{dt}e(t)$ existe, entonces se cumplen ambas para ese instante y además

        \[\frac{d^+}{dt}e(t) = Ae(t) + f(t).\]
    \end{enumerate}
\end{Teorema}

\begin{proof}
    La primera afirmación es trivial. Continuamos entonces con la prueba del tercer punto y la unicidad de soluciones, lo cual probará el segundo punto. Sea $e^* \in D(A^*)$. Para cualquier $e \in E_0$, $t \mapsto \langle e^*, \mathrm{e}^{At}e \rangle$ es diferenciable con derivada $\langle A^*e^*, \mathrm{e}^{At}e \rangle$, pues como $e^*$ es una aplicación lineal y continua, entonces

    \begin{align*}
        \frac{d^+}{dt}\langle e^*, \mathrm{e}^{At}e \rangle &= \left \langle e^*, \frac{d^+}{dt}\mathrm{e}^{At}e \right \rangle\\
        &= \langle e^*, A\mathrm{e}^{At}e \rangle\\
        &= \langle A^* e^*, \mathrm{e}^{At}e \rangle,
    \end{align*}

    donde hemos considerado $e \in D(A)$ y luego, por densidad, extendemos para un $e \in E_0$ cualquiera. Así, por el Teorema fundamental del cálculo:

    \[\langle e^*, \mathrm{e}^{At}e \rangle - \langle e^*, e \rangle = \int_0^t \langle A^* e^*, \mathrm{e}^{As}e \rangle \, ds\]

    para $e \in D(A)$. Como $D(A)$ es denso en $E_0$, y por continuidad del semigrupo $\mathrm{e}^{At}$ y los funcionales $e^*$ y $A^*e^*$, entonces la expresión anterior es válida para $e \in E_0$.\footnotemark
    \footnotetext{Esto nos permite, en el desarrollo que ocurre en las siguientes líneas, poder derivar sin temor a que el dato inicial $e_0$ no esté en el dominio de $A$.}

    Asumimos ahora que
    
    \[e(t) = \mathrm{e}^{A(t-t_0)}e_0 + \int_{t_0}^{t} \mathrm{e}^{A(t-s)}f(s)\, ds\]
    
    y usaremos los resultados anteriores para calcular $\displaystyle\frac{d^+}{dt}\langle e^*, e(t) \rangle$ y comprobar que $e : [t_0, t_1) \rightarrow E_0$ es una solución débil (de acuerdo a la definición \eqref{ch:cauchy:section:noHom:def:soluciones} de solución débil). En efecto:

    \begin{align*}
        \frac{d^+}{dt} \Big \langle e^*, e(t) \Big \rangle &= \frac{d^+}{dt} \left \langle e^*, \mathrm{e}^{A(t-t_0)}e_0 \right \rangle + \frac{d^+}{dt}\left \langle e^*, \int_{t_0}^t \mathrm{e}^{A(t-s)}f(s) \, ds \right \rangle\\
        &=\frac{d^+}{dt} \left( \cancelto{0}{\langle e^*, e_0 \rangle} + \int_0^t \langle A^* e^*, \mathrm{e}^{A(s-t_0)}e_0 \rangle \, ds \right) + \frac{d^+}{dt}\left \langle e^*, \int_{t_0}^t \mathrm{e}^{A(t-s)}f(s) \, ds \right \rangle\\
        &=\langle A^* e^*, \mathrm{e}^{A(t-t_0)}e_0 \rangle + \frac{d^+}{dt}\left \langle e^*, \int_{t_0}^t \mathrm{e}^{A(t-s)}f(s) \, ds \right \rangle.
    \end{align*}

    Vamos a centrarnos ahora en el segundo sumando de la expresión anterior. Construiremos la derivada paso a paso:

    \begin{align*}
        &\left \langle e^*, \int_{t_0}^{t + h}\mathrm{e}^{A(t + h - s)}f(s) \, ds \right \rangle - \left \langle e^*, \int_{t_0}^{t}\mathrm{e}^{A(t - s)}f(s) \, ds \right \rangle\\
        =& \left \langle e^*, \int_{t_0}^{t}\mathrm{e}^{A(t + h - s)}f(s) \, ds + \int_{t}^{t + h}\mathrm{e}^{A(t + h - s)}f(s) \, ds \right \rangle - \left \langle e^*, \int_{t_0}^{t}\mathrm{e}^{A(t - s)}f(s) \, ds \right \rangle\\
        =& \left \langle e^*, \mathrm{e}^{Ah}\int_{t_0}^{t}\mathrm{e}^{A(t - s)}f(s) \, ds + \int_{t}^{t + h}\mathrm{e}^{A(t + h - s)}f(s) \, ds \right \rangle - \left \langle e^*, \int_{t_0}^{t}\mathrm{e}^{A(t - s)}f(s) \, ds \right \rangle\\
        =& \left \langle e^*, \int_{t}^{t + h}\mathrm{e}^{A(t + h - s)}f(s) \, ds \right \rangle + \left \langle e^*, \big(\mathrm{e}^{Ah} - I \big)\int_{t_0}^{t}\mathrm{e}^{A(t - s)}f(s) \, ds \right \rangle\\
        =& \langle e^*, h f(t) \rangle + \left \langle \big(\mathrm{e}^{Ah} - I \big)^* e^*, \int_{t_0}^{t}\mathrm{e}^{A(t - s)}f(s) \, ds \right \rangle + o(h)\\
        =& \langle e^*, h f(t) \rangle + \left \langle \Big(\mathrm{(e^{Ah})^*} - I\Big) e^*, \int_{t_0}^{t}\mathrm{e}^{A(t - s)}f(s) \, ds \right \rangle + o(h)\\
    \end{align*}

    Para continuar, debemos tener en cuenta ahora que $A^*$ es el generador infinitesimal del semigrupo fuertemente continuo $\{T^\odot (t) = \big(T(t)\big)^* |_Y : t \geq 0 \}$ en $L(Y)$, con $Y = \overline{D(A^*)}^{E_0^*}$ \textcolor{red}{(CITAR)}. Luego, como $e^* \in Y$ (pues $e^* \in D(A^*)$), entonces

    \[A^* e^* = \frac{\Big(\mathrm{(e^{Ah})^*} - I\Big) e^*}{h} + o(h) \implies hA^*e^* + o(h) = \Big(\mathrm{(e^{Ah})^*} - I\Big) e^*,\]

    y así

    \begin{align*}
        &\langle e^*, h f(t) \rangle + \left \langle \Big(\mathrm{(e^{Ah})^*} - I\Big) e^*, \int_{t_0}^{t}\mathrm{e}^{A(t - s)}f(s) \, ds \right \rangle + o(h)\\
        =&\langle e^*, h f(t) \rangle + \left \langle hA^* e^*, \int_{t_0}^{t}\mathrm{e}^{A(t - s)}f(s) \, ds \right \rangle + o(h)\\
        =&h \left[ \langle e^*, f(t) \rangle + \left \langle A^* e^*, \int_{t_0}^{t}\mathrm{e}^{A(t - s)}f(s) \, ds \right \rangle \right] + o(h).
    \end{align*}

    En consecuencia:

    \begin{align*}
        &\frac{d^+}{dt} \left \langle e^*, \int_{t_0}^t \mathrm{e}^{A(t-s)}f(s) \, ds\right \rangle\\
        =& \lim_{h \rightarrow 0^+} \frac{1}{h} h \left[ \langle e^*, f(t) \rangle + \left \langle A^* e^*, \int_{t_0}^{t}\mathrm{e}^{A(t - s)}f(s) \, ds \right \rangle \right] + \cancelto{0}{\frac{o(h)}{h}}\\
        =& \langle e^*, f(t) \rangle + \left \langle A^* e^*, \int_{t_0}^{t}\mathrm{e}^{A(t - s)}f(s) \, ds \right \rangle.
    \end{align*}

    Sustituyendo ahora en la expresión de $\frac{d^+}{dt} \Big \langle e^*, e(t) \Big \rangle$:

    \begin{align*}
        \frac{d^+}{dt} \Big \langle e^*, e(t) \Big \rangle &= \langle A^* e^*, \mathrm{e}^{A(t-t_0)}e_0 \rangle + \langle e^*, f(t) \rangle + \left \langle A^* e^*, \int_{t_0}^{t}\mathrm{e}^{A(t - s)}f(s) \, ds \right \rangle\\
        &= \left \langle A^* e^*, \mathrm{e}^{A(t-t_0)}e_0 + \int_{t_0}^{t}\mathrm{e}^{A(t - s)}f(s) \, ds \right \rangle + \langle e^*, f(t) \rangle\\
        &=  \left \langle A^* e^*, e(t) \right \rangle + \langle e^*, f(t) \rangle,
    \end{align*}

    lo que prueba que $e(t)$ es solución débil.

    Continuamos ahora con la prueba de unicidad. Si existen dos soluciones débiles diferentes $e_1(t)$ y $e_2(t)$ de \eqref{ch:cauchy:section:noHom:eqBasica}, entonces denotamos por $u : [t_0, t_1) \rightarrow E_0$ a la diferencia entre ellas. Naturalmente, $u$ es una función continua con $u(t_0) = 0$ y 

    \begin{align*}
        \frac{d^+}{dt}\langle e^*, u(t) \rangle &= \frac{d^+}{dt} \langle e^*, e_1(t) \rangle -  \frac{d^+}{dt} \langle e^*, e_2(t) \rangle\\
        &= \langle A^* e^*, e_1(t) \rangle + \cancel{\langle e^*, f(t) \rangle} - \langle A^* e^*, e_2(t) \rangle - \cancel{\langle e^*, f(t) \rangle}\\
        &= \langle A^* e^*, e_1(t) - e_2(t) \rangle\\
        &= \langle A^* e^*, u(t) \rangle,
    \end{align*}

    para $t_0 \leq t < t_1$ y $e^* \in D(A^*)$.

    Es conveniente trabajar con funciones de clase $C^1$, por lo que definimos

    \[U(t) = \int_{t_0}^t u(s) \, ds.\]

    Luego, tenemos que, integrando a ambos lados de la igualdad $\displaystyle \frac{d^+}{dt}\langle e^*, u(t) \rangle = \langle A^* e^*, u(t) \rangle$:

    \[\langle e^*, u(t) \rangle = \int_{t_0}^t \langle A^* e^*, u(s) \rangle \, ds,\]

    y $\displaystyle \left \langle e^* , \frac{d^+}{dt}U(t) \right \rangle = \langle A^* e^*, U(t) \rangle$. Esto último se deduce de:

    \begin{align*}
        \langle e^*, u(t) \rangle = \int_{t_0}^t \frac{d^+}{dt} \langle e^*, u(s) \rangle \, ds &= \int_{t_0}^t \langle A^* e^*, u(s) \rangle \, ds\\
        &= \left \langle A^* e^*, \int_{t_0}^t u(s) \, ds \right \rangle\\
        &= \langle A^* e^*, U(t) \rangle.
    \end{align*}

    Asimismo, como $U(t) = \int_{t_0}^t u(s) \, ds$, entonces $\displaystyle \frac{d^+}{dt} U(t) = u(t)$. Uniendo ambas expresiones, obtenemos la igualdad deseada.

    Obsérvese ahora que $(\mathrm{e}^{At})^* D(A^*) \subset D(A^*)$ para $t \geq 0$ ya que, para $e^* \in D(A^*)$ y $e \in D(A)$:

    \[ \big \langle (\mathrm{e}^{At})^* e^*, Ae \big \rangle = \big \langle e^*, (\mathrm{e}^{At})Ae \big \rangle = \big \langle e^*, A(\mathrm{e}^{At})e \big \rangle = \big \langle A^*e^*,(\mathrm{e}^{At})e \big \rangle, \]

    por lo que $(\mathrm{e}^{At})^* e^* \in D(A^*).$ Luego, para cualquier $t^* \in (t_0, t_1),$

    \[\left \langle e^*, \mathrm{e}^{A(t^* - t)}\frac{d^+}{dt}U(t)\right \rangle = \langle A^* e^*, \mathrm{e}^{A(t^* - t)}U(t) \rangle.\]

    Armados con todas estas igualdades, tomamos la derivada de $\langle e^*, \mathrm{e}^{A(t^* - t)}U(t) \rangle$:

    \begin{align*}
        \frac{d}{dt} \langle e^*, \mathrm{e}^{A(t^* - t)}U(t) \rangle &= \frac{d}{dt}\Big\langle \big(\mathrm{e}^{A(t^*-t)} \big)^*e^*,U(t)\Big\rangle \\
        &= \left \langle \frac{d}{dt}\big(\mathrm{e}^{A(t^*-t)} \big)^*e^*, U(t) \right \rangle + \left\langle \big(\mathrm{e}^{A(t^*-t)} \big)^*e^*,\frac{d}{dt}U(t)\right\rangle\\
        &= \left \langle - A^*\big(\mathrm{e}^{A(t^*-t)} \big)^*e^*, U(t) \right \rangle + \left\langle \big(\mathrm{e}^{A(t^*-t)} \big)^*e^*, u(t)\right\rangle \\
        &= - \langle A^* e^*, \mathrm{e}^{A(t^* - t)}U(t) \rangle + \langle e^*, \mathrm{e}^{A(t^* - t)} u(t) \rangle\\
        &= - \langle A^* e^*, \mathrm{e}^{A(t^* - t)}U(t) \rangle + \langle A^* e^*, \mathrm{e}^{A(t^* - t)} U(t) \rangle\\
        &= 0
    \end{align*}

    para $t_0 \leq t \leq t^*$.

    Finalmente, tenemos por un lado que $U(t_0) = 0$ y, por otro, que la función $\langle e^*, \mathrm{e}^{A(t^* - t)}U(t) \rangle$ es contante para $t_0 \leq t \leq t^*$. Como $U(t_0) = 0$, entonces dicha función debe ser nula también en $t = t_0$. Así, para $t = t^*$:

    \[\langle e^*, \mathrm{e}^{A(t^* - t^*)}U(t^*) \rangle = \langle e^*, U(t^*) \rangle = 0 \ ; \ t^* \in (t_0, t_1).\]

    Como lo anterior es cierto para todo $e^* \in D(A^*)$, entonces $U(t^*) =0$. Esto lo sabemos porque $A$ es cerrado (pues es el generador de un semigrupo fuertemente continuo e invocando el teorema de Hille-Yosida \eqref{ch:semigrupos:th:HilleYosida}) se sigue que esto implica que es cerrado) y densamente definido. Ambas cosas implican, por el Lema \eqref{ch:cauchy:section:noHom:lema:dominioAdjTotal}, que $D(A^*)$ es total, por lo que $U(t^*) = 0$.
    
    Así, $u(s) = 0$ para todo $s \in (t_0, t_1)$. Esto prueba entonces la unicidad de soluciones y, junto con el punto tercero, también prueba el segundo punto.

    Concluimos probando el cuarto y último punto. Sea $e(\cdot)$ una solución débil dada por \eqref{ch:cauchy:section:noHom:th:solDebilYFuerteRels:2}, entonces para $t_0 \leq t < t + h \leq t$:

    \begin{align*}
        &\frac{e(t+h) - e(t)}{h}\\
        &= \frac{1}{h} \left[ \mathrm{e}^{A(t+h-t_0)}e_0 + \int_{t_0}^{t+h}\mathrm{e}^{A(t+h-s)} f(s) \, ds - \mathrm{e}^{A(t-t_0)}e_0 - \int_{t_0}^{t}\mathrm{e}^{A(t-s)} f(s) \, ds \right]\\
        &= \frac{1}{h} \left[ \mathrm{e}^{Ah}\mathrm{e}^{A(t-t_0)}e_0 + \mathrm{e}^{Ah}\int_{t_0}^{t}\mathrm{e}^{A(t-s)} f(s) \, ds + \int_{t}^{t+h}\mathrm{e}^{A(t+h-s)} f(s) \, ds - \mathrm{e}^{A(t-t_0)}e_0 - \int_{t_0}^{t}\mathrm{e}^{A(t-s)} f(s) \, ds \right]\\
        &= \frac{1}{h}\left[\int_{t}^{t+h}\mathrm{e}^{A(t+h-s)} f(s) \, ds + (\mathrm{e}^{Ah} - I) \left( \mathrm{e}^{A(t-t_0)}e_0 + \int_{t_0}^{t}\mathrm{e}^{A(t-s)} f(s) \, ds \right) \right]\\
        &= \frac{1}{h} \left[\int_{t}^{t+h}\mathrm{e}^{A(t+h-s)} f(s) \, ds + (\mathrm{e}^{Ah} - I)e(t) \right]\\
        &= \frac{1}{h} \int_{t}^{t+h}\mathrm{e}^{A(t+h-s)} f(s) \, ds + \frac{1}{h}(\mathrm{e}^{Ah} - I)e(t).
    \end{align*}

    Luego, el término

    \[\frac{1}{h} \int_{t}^{t+h}\mathrm{e}^{A(t+h-s)} f(s) \, ds\]

    converge a $f(t^+) = f(t)$ (dicha igualdad se da por la continuidad de $f$) cuando $h \rightarrow 0^+$. Luego, si cualquiera de los otros dos términos converge, entonces el que resta también debe hacerlo, probando así el cuarto y último punto del Teorema.
\end{proof}

A continuación, damos condiciones simples que aseguran la diferenciabilidad de una solución débil.\\

\begin{Teorema}{(Condiciones de diferenciabilidad para soluciones débiles)}
    \label{ch:cauchy:section:noHom:th:conDiffSolDebil}

    Asúmanse $A$ y $f$ como antes, y sea $e : [t_0, t_1) \rightarrow E_0$ una solución débil de \eqref{ch:cauchy:section:noHom:eqBasica}. Si $e_0 \in D(A)$ y bien
    
    \begin{enumerate}
    \item $f(t) \in D(A)$ con $t \mapsto Af(t) \in E_0$ continua por la derecha en $[t_0, t_1)$, o bien

    \item $\displaystyle \frac{d^+}{dt}f(t) = \dot{f}(t)$ existe y es continua por la derecha en $[t_0, t_1)$
    \end{enumerate}

    entonces $\displaystyle \frac{d^+}{dt}e(t)$ existe, $e(t) \in D(A)$ y $e : [t_0, t_1) \rightarrow E_0$ es una solución fuerte.
\end{Teorema}

\begin{proof}
    Sea $\displaystyle u(t) = \int_{t_0}^t \mathrm{e}^{A(t-s)}f(s) \, ds$, luego $e(t) = \mathrm{e}^{A(t-t_0)}e_0 + u(t)$. Como $e_0 \in D(A)$, entonces $t \mapsto \mathrm{e}^{A(t-t_0)}e_0$ es de clase $C^1$. Sea entonces $t_0 \leq t < t + h < t_1$, se tiene que:

    \begin{align*}
        \frac{u(t+h) - u(t)}{h} &= \frac{1}{h}\int_{t_0}^{t+h} \mathrm{e}^{A(t + h -s)}f(s) \, ds - \frac{1}{h}\int_{t_0}^t \mathrm{e}^{A(t-s)}f(s) \, ds\\
        &= \frac{1}{h}\mathrm{e}^{Ah}\int_{t_0}^{t} \mathrm{e}^{A(t-s)}f(s) \, ds + \frac{1}{h}\int_{t}^{t+h} \mathrm{e}^{A(t + h -s)}f(s) \, ds - \frac{1}{h}\int_{t_0}^t \mathrm{e}^{A(t-s)}f(s) \, ds\\
        &= \frac{1}{h}\int_{t}^{t+h} \mathrm{e}^{A(t + h -s)}f(s) \, ds + \int_{t_0}^t \mathrm{e}^{A(t-s)}\frac{\mathrm{e}^{Ah} - I}{h}f(s) \, ds
    \end{align*}

    Por otro lado, también podemos desarrollar la expresión anterior como:

    \begin{align*}
        \frac{u(t+h) - u(t)}{h} &= \frac{1}{h}\int_{t_0}^{t+h} \mathrm{e}^{A(t + h -s)}f(s) \, ds - \frac{1}{h}\int_{t_0}^t \mathrm{e}^{A(t-s)}f(s) \, ds\\
        &= \frac{1}{h}\int_{t_0}^{t} \mathrm{e}^{A(t + h -s)}f(s) \, ds + \frac{1}{h}\int_{t}^{t+h} \mathrm{e}^{A(t + h -s)}f(s) \, ds - \frac{1}{h}\int_{t_0}^t \mathrm{e}^{A(t-s)}f(s) \, ds\\
        \text{(Ca. de var.) }&= \frac{1}{h}\int_{t_0 - h}^{t-h} \mathrm{e}^{A(t-x)}f(x + h) \, dx + \frac{1}{h}\int_{t}^{t+h} \mathrm{e}^{A(t + h -s)}f(s) \, ds - \frac{1}{h}\int_{t_0}^t \mathrm{e}^{A(t-s)}f(s) \, ds\\
        &= \frac{1}{h}\Bigg[ \int_{t_0 - h}^{t_0} \mathrm{e}^{A(t-x)}f(x + h) \, dx + \int_{t_0}^{t} \mathrm{e}^{A(t-x)}f(x + h) \, dx\\
        &- \int_{t - h}^{t} \mathrm{e}^{A(t-x)}f(x + h) \, dx\Bigg ] + \frac{1}{h}\int_{t}^{t+h} \mathrm{e}^{A(t + h -s)}f(s) \, ds - \frac{1}{h}\int_{t_0}^t \mathrm{e}^{A(t-s)}f(s) \, ds\\
        \text{(Ca. de var.) }&=\frac{1}{h}\Bigg[ \int_{t_0}^{t_0 + h} \mathrm{e}^{A(t+h-s)}f(s) \, ds + \int_{t_0}^{t} \mathrm{e}^{A(t-x)}f(x + h) \, dx\\
        &- \cancel{\int_{t}^{t+h} \mathrm{e}^{A(t+h-s)}f(s) \, ds}\Bigg ] + \cancel{\frac{1}{h}\int_{t}^{t+h} \mathrm{e}^{A(t + h -s)}f(s) \, ds} - \frac{1}{h}\int_{t_0}^t \mathrm{e}^{A(t-s)}f(s) \, ds\\
        &= \frac{1}{h} \int_{t_0}^{t_0 + h} \mathrm{e}^{A(t+h-s)}f(s) \, ds + \frac{1}{h}\int_{t_0}^{t} \mathrm{e}^{A(t-x)}f(x + h) \, dx - \frac{1}{h}\int_{t_0}^t \mathrm{e}^{A(t-s)}f(s) \, ds\\
        &= \frac{1}{h} \int_{t_0}^{t_0 + h} \mathrm{e}^{A(t+h-s)}f(s) \, ds + \int_{t_0}^t \mathrm{e}^{A(t-s)}\frac{f(s+h) - f(s)}{h} \, ds.
    \end{align*}

    Teniendo ambas expresiones, ahora tomamos el límite cuando $h \rightarrow 0^+$ de la primera de ellas e, imponiendo la primera condición, obtenemos:

    \[\frac{d^+}{dt}u(t) = f(t) + \int_{t_0}^t \mathrm{e}^{A(t-s)}Af(s) \, ds.\]

    Es decir, $\displaystyle \frac{d^+}{dt}u(t)$ existe y $\displaystyle t \mapsto \frac{d^+}{dt}u(t)$ es continua.

    Repetimos ahora con la segunda expresión. Tomamos el límite e imponemos la segunda condición, obteniendo así:

    \[\frac{d^+}{dt}u(t) = \mathrm{e}^{A(t-t_0)}f(t_0) + \int_{t_0}^t \mathrm{e}^{A(t-s)}\dot{f}(s) \, ds,\]

    que, de nuevo, implica que $\displaystyle \frac{d^+}{dt}u(t)$ existe y que $\displaystyle t \mapsto \frac{d^+}{dt}u(t)$ es continua.

    Finalmente, por lo dicho al comienzo (entiéndase: que $t \mapsto \mathrm{e}^{A(t-t_0)}e_0$ es de clase $C^1$ por ser $e_0$ un elemento de $D(A)$), se tiene que $\displaystyle \frac{d^+}{dt}e(t)$ existe (pues $\displaystyle \frac{d^+}{dt}e(t) = \frac{d^+}{dt}\Big(\mathrm{e}^{A(t_0 - t)}e_0 \Big) + \frac{d^+}{dt}u(t)$) y $\displaystyle t \mapsto \frac{d^+}{dt}e(t)$ es continua. Así, por el cuarto punto del Teorema \eqref{ch:cauchy:section:noHom:th:solDebilYFuerteRels}

    \[\frac{d^+}{dt}e(t) = Ae(t) + f(t) \ ; \ e(t) \in D(A), \]

    y $e: [t_0, t_1) \rightarrow E_0$ es una solución fuerte.
\end{proof}

Esto concluye entonces nuestro estudio sobre la existencia, unicidad y manipulación de soluciones (tanto débiles como fuertes) del problema de Cauchy abstracto no homogéneo. Pasemos ahora a estudiar el problema abstracto de Cauchy semilineal (en concreto, el caso hiperbólico), apoyándonos en los resultados obtenidos en esta sección.

\section{Problema de Cauchy semilineal, caso hiperbólico}
\label{ch:cauchy:section:semilineal}

Similar a la sección anterior, en esta sección consideraremos el problema de valor inicial \textcolor{red}{ARREGLAR EL LABEL}

\begin{align}
    \label{ch:cauchy:section:semilineal:eqBasica}
    \frac{de}{dt} &= Ae + f(t,e)\\
    e(t_0) &= e_0 \in E_0,
\end{align}

donde $A$ es el generador de un semigrupo fuertemente continuo, $f$ es una función continua que está definida en un subconjunto $U$ de $\mathbb{R} \times E_0$ y toma valores en $E_0$. Además, $(t_0, e_0) \in U$.

Siguiendo la misma estructura de la sección anterior, empezamos por definir los tipos de soluciones que puede tener el problema \eqref{ch:cauchy:section:semilineal:eqBasica}\\

\begin{Definicion}{(Soluciones al problema de Cauchy semilineal, caso hiperbólico)}
    \label{ch:cauchy:section:semilineal:def:soluciones}

    \begin{enumerate}[label=\alph*)]
        \item Una \textit{solución fuerte} de \eqref{ch:cauchy:section:semilineal:eqBasica} en $[t_0, t_1)$ es una función continua $e : [t_0, t_1) \rightarrow E_0$ tal que $e(t_0) = e_0$, $e : (t_0, t_1) \rightarrow E_0$ es continuamente diferenciable, $(t, e(t)) \in U$, $e(t) \in D(A)$ y se verifica \eqref{ch:cauchy:section:semilineal:eqBasica} para $t \in (t_0, t_1)$. 

        \item Una \textit{solución débil} de \eqref{ch:cauchy:section:semilineal:eqBasica} en $[t_0, t_1)$ es una función continua $e : [t_0, t_1) \rightarrow E_0$ tal que $e(t_0) = e_0$, $(t, e(t)) \in U$, para todo $e^* \in D(A^*)$ se tiene que $t \mapsto \langle e^*, e(t) \rangle$ es diferenciable y 

        \[\frac{d}{dt}\langle e^*, e(t) \rangle = \langle A^* e^*, e(t) \rangle + \langle e^*, f(t, e(t)) \rangle, \quad t_0 < t < t_1.\]
    \end{enumerate}
\end{Definicion}

Dada la definición anterior, enunciamos un Teorema análogo al Teorema \eqref{ch:cauchy:section:noHom:th:solDebilYFuerteRels}. Es decir, un Teorema que nos dará información sobre la forma de las soluciones débiles, su relación con las soluciones fuertes y algunas formas de manipularlas.\\

\begin{Teorema}{(Manipulación de soluciones débiles y relaciones entre soluciones débiles y fuertes)}
    \label{ch:cauchy:section:semilineal:th:solDebilYFuerteRels}

    \begin{enumerate}
        \item Si $e : [t_0, t_1) \rightarrow E_0$ es una solución fuerte de \eqref{ch:cauchy:section:semilineal:eqBasica}, entonces es también una solución débil de \eqref{ch:cauchy:section:semilineal:eqBasica}.

        \item Una solución débil $e : [t_0, t_1) \rightarrow E_0$ de \eqref{ch:cauchy:section:semilineal:eqBasica} es también una solución fuerte si y sólo si es continuamente diferenciable en $(t_0, t_1)$ si y sólo si $e(t) \in D(A)$ con $t \mapsto Ae(t)$ es continua en $(t_0, t_1)$.

        \item Si $e : [t_0, t_1) \rightarrow E_0$ es una solución débil de \eqref{ch:cauchy:section:semilineal:eqBasica}, entonces

        \begin{align}
            \label{ch:cauchy:section:semilineal:th:solDebilYFuerteRels:2}
            e(t) = \mathrm{e}^{A(t-t_0)}e_0 + \int_{t_0}^{t}\mathrm{e}^{A(t-s)}f(s, e(s)) \, ds, \quad t_0 \leq t \leq t_1
        \end{align}

        \item Si $e : [t_0, t_1) \rightarrow E_0$ es continua con $(t,e(t)) \in U, t_0 \leq t < t_1$, y satisface \eqref{ch:cauchy:section:semilineal:eqBasica}, entonces $e : [t_0, t_1) \rightarrow E_0$ es una solución débil de \eqref{ch:cauchy:section:semilineal:eqBasica}.
    \end{enumerate}
\end{Teorema}

\begin{proof}
    El razonamiento detrás de esta demostración se basa en reducir el problema semilineal al marco de trabajo del problema lineal estudiado previamente. En la sección anterior, el Teorema \eqref{ch:cauchy:section:noHom:th:solDebilYFuerteRels} establece las propiedades fundamentales para el problema de Cauchy lineal no homogéneo, cuya estructura es de la forma:
    
    \[\frac{d}{dt}e(t) = Ae(t) + f(t)\]
    
    donde el término no homogéneo $f(t)$ depende únicamente de la variable temporal $t$. Por otro lado, para el problema de Cauchy semilineal \eqref{ch:cauchy:section:semilineal:eqBasica}, la ecuación diferencial involucra una no linealidad que depende tanto del tiempo como del estado del sistema:
    
    \[\frac{d}{dt}e(t) = Ae(t) + f(t, e(t)).\]
    
    El argumento central de la demostración consiste en fijar una función continua $e: [t_0, t_1) \to E_0$ y definir una nueva función $g: [t_0, t_1) \to E_0$ dada por $g(t) = f(t, e(t))$. Dado que, por hipótesis, la aplicación $t \mapsto f(t, e(t))$ es continua, la función $g(t)$ resulta ser simplemente un término fuente continuo que depende exclusivamente de $t$. Sustituyendo esta nueva función, la ecuación diferencial semilineal se reescribe como:
    
    \[\frac{d}{dt}e(t) = Ae(t) + g(t)\]
    
    Al realizar esta identificación, el problema semilineal adquiere exactamente la misma estructura que el problema lineal. Puesto que $g(t)$ satisface la condición de continuidad exigida sobre el término independiente en la sección anterior, todas las afirmaciones del Teorema \eqref{ch:cauchy:section:noHom:th:solDebilYFuerteRels} se heredan de manera inmediata. Esto incluye la equivalencia entre soluciones débiles y fuertes (bajo las condiciones de regularidad adecuadas) y la representación integral de la solución expresada en la ecuación \eqref{ch:cauchy:section:semilineal:th:solDebilYFuerteRels:2}, obtenida simplemente al reemplazar el término lineal por $f(s, e(s))$ en la integral.
\end{proof}

\textcolor{red}{CONTINUAR????...}